\documentclass[12pt, a4paper]{article}

\usepackage[utf8]{inputenc}
\usepackage[T1]{fontenc}
\usepackage[margin=2.5cm]{geometry}
\usepackage{hyperref}
\usepackage{listings}
\usepackage{xcolor}
\usepackage{booktabs}
\usepackage{longtable}
\usepackage{array}
\usepackage{enumitem}
\usepackage{titlesec}
\usepackage{fancyhdr}
\usepackage{graphicx}
\usepackage{amsmath}

% ── Colors ──────────────────────────────────────────────────────────────────
\definecolor{dxcpurple}{RGB}{102, 0, 153}
\definecolor{codebg}{RGB}{245, 245, 245}
\definecolor{codeframe}{RGB}{200, 200, 200}
\definecolor{keywordcolor}{RGB}{0, 0, 180}
\definecolor{stringcolor}{RGB}{163, 21, 21}
\definecolor{commentcolor}{RGB}{0, 128, 0}

% ── Code style ──────────────────────────────────────────────────────────────
\lstdefinestyle{python}{
    language=Python,
    backgroundcolor=\color{codebg},
    frame=single,
    rulecolor=\color{codeframe},
    basicstyle=\ttfamily\small,
    keywordstyle=\color{keywordcolor}\bfseries,
    stringstyle=\color{stringcolor},
    commentstyle=\color{commentcolor}\itshape,
    numbers=left,
    numberstyle=\tiny\color{gray},
    numbersep=8pt,
    breaklines=true,
    breakatwhitespace=true,
    tabsize=4,
    showstringspaces=false,
    captionpos=b,
}

\lstdefinestyle{bash}{
    language=bash,
    backgroundcolor=\color{codebg},
    frame=single,
    rulecolor=\color{codeframe},
    basicstyle=\ttfamily\small,
    keywordstyle=\color{keywordcolor}\bfseries,
    stringstyle=\color{stringcolor},
    commentstyle=\color{commentcolor}\itshape,
    breaklines=true,
    showstringspaces=false,
}

\lstdefinestyle{json}{
    backgroundcolor=\color{codebg},
    frame=single,
    rulecolor=\color{codeframe},
    basicstyle=\ttfamily\small,
    breaklines=true,
    showstringspaces=false,
}

% ── Header / Footer ─────────────────────────────────────────────────────────
\pagestyle{fancy}
\fancyhf{}
\fancyhead[L]{\textcolor{dxcpurple}{\textbf{DXC Technology}}}
\fancyhead[R]{\textcolor{dxcpurple}{Jira API Integration}}
\fancyfoot[C]{\thepage}
\renewcommand{\headrulewidth}{0.4pt}
\renewcommand{\headrule}{\hbox to\headwidth{\color{dxcpurple}\leaders\hrule height \headrulewidth\hfill}}

% ── Section colors ───────────────────────────────────────────────────────────
\titleformat{\section}{\large\bfseries\color{dxcpurple}}{}{0em}{}[\titlerule]
\titleformat{\subsection}{\normalsize\bfseries\color{dxcpurple}}{}{0em}{}

% ── Hyperref ────────────────────────────────────────────────────────────────
\hypersetup{
    colorlinks=true,
    linkcolor=dxcpurple,
    urlcolor=dxcpurple,
    citecolor=dxcpurple,
}

% ────────────────────────────────────────────────────────────────────────────
\begin{document}

% ── Title page ──────────────────────────────────────────────────────────────
\begin{titlepage}
    \centering
    \vspace*{2cm}
    {\color{dxcpurple}\rule{\linewidth}{1.5pt}}\\[0.5cm]
    {\Huge\bfseries\color{dxcpurple} Jira API Integration}\\[0.3cm]
    {\Large Technical Procedure Document}\\[0.3cm]
    {\color{dxcpurple}\rule{\linewidth}{1.5pt}}\\[2cm]

    {\large\textbf{Project:} IRP AUTO — DXC GraphTalk Service Desk Reporting System}\\[0.5cm]
    {\large\textbf{Instance:} \texttt{dxc-insurance-delivery.atlassian.net}}\\[0.5cm]
    {\large\textbf{Project Key:} \texttt{IRPAUTO}}\\[2cm]

    \begin{tabular}{ll}
        \textbf{Author}   & El Adnani Ali \\[0.3cm]
        \textbf{Date}     & April 2, 2026 \\[0.3cm]
        \textbf{Version}  & 1.0 \\[0.3cm]
        \textbf{Status}   & Production \\
    \end{tabular}

    \vfill
    {\color{dxcpurple}\rule{\linewidth}{0.4pt}}\\[0.2cm]
    {\small DXC Technology — IRP AUTO Service Desk Analytics Platform}
\end{titlepage}

\tableofcontents
\newpage

% ────────────────────────────────────────────────────────────────────────────
\section{Overview}

This document describes the end-to-end procedure for connecting the DXC IRP AUTO reporting
dashboard to Atlassian Jira Cloud via the REST API v3. Prior to this integration, issue data
was obtained by manually exporting an Excel file from Jira and uploading it through the
application. This process was replaced by a direct programmatic connection that allows
on-demand and incremental synchronisation without any manual file handling.

\subsection{Objectives}

\begin{itemize}
    \item Eliminate the manual Excel export/upload workflow.
    \item Provide full parity with all columns previously available in the Excel extract.
    \item Support both \textbf{incremental sync} (issues updated since a given date) and
          \textbf{full reload} (all issues from the beginning of the project).
    \item Expose a user-friendly sync interface inside the existing Streamlit dashboard.
\end{itemize}

\subsection{Architecture Summary}

The integration follows a straightforward \textbf{Extract–Transform–Load (ETL)} pattern:

\begin{enumerate}
    \item \textbf{Extract} — Paginate through the Jira REST API v3 \texttt{/search/jql}
          endpoint using a JQL query.
    \item \textbf{Transform} — Map raw Jira field values to the application's column schema,
          parse dates, resolve numeric types, and compute derived columns.
    \item \textbf{Load} — Write the resulting DataFrame into the MySQL database (Aiven Cloud)
          using either an upsert or a full replace strategy.
\end{enumerate}

% ────────────────────────────────────────────────────────────────────────────
\section{Prerequisites}

\subsection{Jira API Token}

Authentication uses an Atlassian \textbf{API Token} (not a password). To generate one:

\begin{enumerate}
    \item Log in to \href{https://id.atlassian.com/manage-profile/security/api-tokens}
          {\texttt{id.atlassian.com}} with your DXC account.
    \item Click \textit{Create API token}, give it a name, and copy the generated value.
    \item Store it in the application's \texttt{.env} file (see Section~\ref{sec:env}).
\end{enumerate}

\subsection{Python Dependency}

The only additional library required beyond what was already installed is \texttt{requests}:

\begin{lstlisting}[style=bash]
pip install requests
\end{lstlisting}

\subsection{Environment Variables}
\label{sec:env}

Three new variables are added to the \texttt{.env} file:

\begin{lstlisting}[style=bash]
JIRA_URL=https://dxc-insurance-delivery.atlassian.net
JIRA_EMAIL=ali.eladnani@dxc.com
JIRA_API_TOKEN=<your-api-token>
\end{lstlisting}

These are read at runtime by both \texttt{etl.py} and \texttt{app.py} via \texttt{python-dotenv}.

% ────────────────────────────────────────────────────────────────────────────
\section{API Discovery}

Before implementing the integration, the Jira REST API was explored interactively to verify
connectivity, discover available projects, and map every field used by the existing Excel
extract to its corresponding Jira API field identifier.

\subsection{Authentication Test}

\begin{lstlisting}[style=bash, caption={Verify API credentials}]
curl -u "email:api_token" \
  "https://dxc-insurance-delivery.atlassian.net/rest/api/3/myself"
\end{lstlisting}

A successful response returns the authenticated user's profile, confirming that Basic Auth
with an API token works correctly against this Jira Cloud instance.

\subsection{Project Discovery}

\begin{lstlisting}[style=bash, caption={List accessible projects}]
curl -u "email:api_token" \
  "https://dxc-insurance-delivery.atlassian.net/rest/api/3/project/search"
\end{lstlisting}

The IRPAUTO project (\textit{IRP AUTO — DXC GraphTalk Service Desk}) was confirmed accessible
among five projects: DSP, GLPS, IRPAUTO, SRSI, and TEST.

\subsection{Field Mapping Discovery}

The complete list of available fields was retrieved from:

\begin{lstlisting}[style=bash, caption={Retrieve all field definitions}]
curl -u "email:api_token" \
  "https://dxc-insurance-delivery.atlassian.net/rest/api/3/field"
\end{lstlisting}

Individual issues (including a known closed ticket and a scope ticket) were inspected
with \texttt{?fields=*all} to identify which custom fields carry meaningful values for IRPAUTO.

\subsection{Issue Count}

The approximate total was retrieved via the count endpoint:

\begin{lstlisting}[style=bash, caption={Approximate issue count}]
POST /rest/api/3/search/approximate-count
{ "jql": "project = IRPAUTO" }
# Response: { "count": 55244 }
\end{lstlisting}

% ────────────────────────────────────────────────────────────────────────────
\section{Field Mapping}

\subsection{Understanding Jira Field Types}

Jira has two categories of fields, which look very different in the API response:

\begin{itemize}
    \item \textbf{Standard fields} — built into every Jira instance. They have simple names
          like \texttt{assignee}, \texttt{status}, \texttt{priority}. In the API response
          they are objects containing sub-properties. For example, the \texttt{assignee}
          field returns a JSON object like:
          \begin{lstlisting}[style=json]
"assignee": {
    "displayName": "El Adnani Ali",
    "emailAddress": "ali.eladnani@dxc.com",
    "accountId": "712020:2dce71f4-..."
}
          \end{lstlisting}
          Our ETL extracts the \texttt{displayName} sub-property. There is no separate field
          called \texttt{assignee.displayName} — it is simply the \textit{name of the person}
          stored inside the \texttt{assignee} object.

    \item \textbf{Custom fields} — created specifically for this Jira instance by an
          administrator. They have auto-generated IDs like \texttt{customfield\_10492}.
          In the Jira UI they appear under their human-readable label (e.g.\ ``Qualification
          date''), but in the API they are always referenced by their numeric ID.
\end{itemize}

\subsection{How to Find a Custom Field ID in Jira}

There are three ways to look up the ID of any custom field:

\begin{enumerate}
    \item \textbf{Jira Admin panel} (requires admin access):\\
          Navigate to \textit{Settings} $\to$ \textit{Issues} $\to$ \textit{Custom fields}.
          Click on any field name. The URL in your browser will contain the field ID,
          e.g.\ \texttt{.../customFieldId=10492}.

    \item \textbf{REST API field list} (no admin required):\\
          Call \texttt{GET /rest/api/3/field} — this returns the full list of all fields
          with their IDs and human-readable names. Filter the result for the field name
          you are looking for.

    \item \textbf{Issue view in browser}:\\
          Open any issue in Jira, right-click on a field label, and inspect the element.
          The HTML often contains the field ID in a \texttt{data-field-id} attribute.
\end{enumerate}

\subsection{Complete Field Reference Table}

Table~\ref{tab:fieldmap} maps every Excel column to the corresponding Jira field. The
\textit{Jira UI Label} column shows the name as it appears on the issue screen in the
browser. The \textit{API Field ID} column is what the code uses to request and read the
field from the REST API.

\begin{longtable}{>{\small}p{3.6cm} >{\ttfamily\small}p{3.2cm} >{\small}p{3.0cm} >{\ttfamily\small}p{3.0cm}}
    \toprule
    \normalfont\textbf{Excel Column} &
    \normalfont\textbf{DB Column} &
    \normalfont\textbf{Jira UI Label} &
    \normalfont\textbf{API Field ID} \\
    \midrule
    \endhead
    \bottomrule
    \endfoot
    \caption{Complete field mapping: Excel $\to$ DB $\to$ Jira UI $\to$ API ID}
    \label{tab:fieldmap}
    \endlastfoot

    Key                        & key                        & Key                          & issuekey \\
    Issue Type                 & issue\_type                & Issue Type                   & issuetype \\
    Summary                    & summary                    & Summary                      & summary \\
    Assignee                   & assignee                   & Assignee                     & assignee \\
    Reporter                   & reporter                   & Reporter                     & reporter \\
    Priority                   & priority                   & Priority                     & priority \\
    Status                     & status                     & Status                       & status \\
    Created                    & created                    & Created                      & created \\
    Resolved                   & resolved                   & Resolved                     & resolutiondate \\
    Due date                   & due\_date                  & Due date                     & duedate \\
    Fix versions               & fix\_versions              & Fix versions                 & fixVersions \\
    Project                    & project                    & Project                      & project \\
    Closure Date               & closure\_date              & Closure Date                 & customfield\_10297 \\
    Qualification date         & qualification\_date        & Qualification date           & customfield\_10492 \\
    Sending Date               & sending\_date              & Sending Date                 & customfield\_10369 \\
    Last Comment Date/Time     & last\_comment\_datetime    & Last Comment Date/Time       & customfield\_10248 \\
    Start Date \& Time         & start\_date\_time          & Start Date \& Time           & customfield\_10317 \\
    Story Points               & story\_points              & Story Points (numeric)       & customfield\_10049 \\
    Story Points (IRPAUTO)     & story\_points              & Story Points (dropdown)*     & customfield\_10520 \\
    Organizations              & organizations              & Organizations                & customfield\_10002 \\
    Customer Reference Id      & customer\_reference\_id    & Customer Reference Id        & customfield\_10127 \\
    Environment Type           & environment\_type          & Environment Type             & customfield\_10187 \\
    Environment                & environment                & Environment                  & customfield\_10238 \\
    Component                  & component                  & Component                    & customfield\_10171 \\
    Current Tier               & current\_tier              & Current Tier                 & customfield\_10173 \\
    Reproductibility           & reproductibility           & Reproductibility             & customfield\_10242 \\
    Resolutions                & resolutions                & Resolutions                  & customfield\_10301 \\
    Root Cause Description     & root\_cause\_description   & Root Cause Description       & customfield\_10310 \\
    Impact Description         & impact\_description        & Impact Description           & customfield\_10312 \\
    Resolution Description     & resolution\_description    & Resolution Description       & customfield\_10311 \\
    Client Priority            & client\_priority           & Client Priority              & customfield\_10303 \\
    Target Release             & target\_release            & Target Release               & customfield\_10179 \\
    Count Ready Acceptance     & count\_ready\_acceptance   & Count of Ready for acceptance & customfield\_10327 \\
    Count Ready Staging        & count\_ready\_staging      & Count of ready for staging   & customfield\_10321 \\
    Product Line               & product\_line              & Product Line                 & customfield\_10229 \\
    Occurrences                & occurrences                & Occurrences                  & customfield\_10381 \\
    SLA Justified              & sla\_justified             & SLA Justified                & customfield\_10429 \\
    OUT OF SLA Reason          & out\_of\_sla\_reason       & OUT OF SLA Reason            & customfield\_10446 \\
    Root Cause Origin          & root\_cause\_origin        & Root Cause Origin            & customfield\_10452 \\
    Nombre d'UO                & nombre\_uo                 & Nombre d'UO                  & customfield\_10447 \\
    Service request types      & service\_request\_types    & Service request types        & customfield\_10165 \\
    Enhancement request type   & enhancement\_request\_type & Enhancement request type     & customfield\_10382 \\
    Resolution Owner           & resolution\_owner          & Resolution Owner             & customfield\_11038 \\
    Expected Sprint            & expected\_sprint           & Expected Sprint              & customfield\_10315 \\
    Original Expected Sprint   & original\_expected\_sprint & Original Expected Sprint     & customfield\_10642 \\
    Baseline / Standard Issue  & baseline\_standard\_issue  & Baseline / Standard Issue    & customfield\_10305 \\
    KPI - Time to Solve        & kpi\_time\_to\_solve\_minutes      & KPI - Time to Solve$^\dagger$  & customfield\_10186 \\
    KPI - Time Initial Response & kpi\_time\_initial\_response\_minutes & KPI - Time Initial Response$^\dagger$ & customfield\_10185 \\
    KPI - Time to Analyse      & kpi\_time\_to\_analyse\_minutes    & KPI - Time to Analyse$^\dagger$ & customfield\_10485 \\
    KPI - Time to Assist       & kpi\_time\_to\_assist\_minutes     & KPI - Time to Assist$^\dagger$ & customfield\_10428 \\
    KPI - Time to Estimate     & kpi\_time\_to\_estimate\_minutes   & KPI - Time to Estimate$^\dagger$ & customfield\_10293 \\

\end{longtable}

\noindent
* \textbf{Story Points — dual field (IRPAUTO specific):} IRPAUTO uses two different fields
for story points. \texttt{customfield\_10049} is the standard Jira numeric SP field, but it
is empty for \textasciitilde82\% of IRPAUTO tickets. The actual values are stored in
\texttt{customfield\_10520}, which is a dropdown field whose options are the point values
(\texttt{1}, \texttt{2}, \texttt{3}, \texttt{5}, \texttt{8}, \texttt{12}, \texttt{20}).
The ETL reads \texttt{customfield\_10049} first and falls back to \texttt{customfield\_10520}
when null.

\vspace{0.4cm}
\noindent
$\dagger$ \textbf{KPI SLA fields:} In the Jira UI these fields show durations like
\texttt{998h 31m}. However the REST API returns them as complex SLA cycle objects
(start timestamp, stop timestamp, breach status, remaining time) — not a pre-formatted
string. The Excel export pre-processes these into the human-readable format. The ETL
currently stores these as \texttt{NULL}; parsing the raw SLA objects is deferred to a
future iteration.

% ────────────────────────────────────────────────────────────────────────────
\section{Implementation}

\subsection{ETL Function: \texttt{fetch\_jira\_issues}}

The core of the integration is the \texttt{fetch\_jira\_issues} function added to
\texttt{etl.py}. Its full signature is:

\begin{lstlisting}[style=python, caption={Function signature}]
def fetch_jira_issues(
    engine,
    jira_url: str,
    email: str,
    api_token: str,
    project: str = "IRPAUTO",
    mode: str = "upsert",       # "upsert" | "replace"
    updated_since: str = None,  # ISO date "YYYY-MM-DD" for incremental
    progress_cb = None,         # callable(done, total) for UI progress
) -> dict:  # {"fetched": N, "loaded": N}
\end{lstlisting}

\subsection{Step 1 — Approximate Count}

Before fetching, the function queries the approximate count endpoint to enable an accurate
progress bar in the UI:

\begin{lstlisting}[style=python, caption={Pre-fetch count for progress bar}]
resp = requests.post(
    f"{base}/rest/api/3/search/approximate-count",
    auth=auth, headers=headers,
    json={"jql": jql}, timeout=15,
)
total_hint = resp.json().get("count", 0)
\end{lstlisting}

\subsection{Step 2 — JQL Construction}

The JQL query is built dynamically based on whether the sync is incremental or full:

\begin{lstlisting}[style=python, caption={JQL construction}]
jql = f"project = {project}"
if updated_since:
    jql += f' AND updated >= "{updated_since}"'
jql += " ORDER BY updated ASC"
\end{lstlisting}

\subsection{Step 3 — Pagination}

The new \texttt{/rest/api/3/search/jql} endpoint (v3, replacing the deprecated
\texttt{/search}) uses cursor-based pagination via a \texttt{nextPageToken}:

\begin{lstlisting}[style=python, caption={Cursor-based pagination loop}]
next_page_token = None
while True:
    payload = {
        "jql":        jql,
        "maxResults": 100,
        "fields":     list(_JIRA_FIELD_MAP.keys()),
    }
    if next_page_token:
        payload["nextPageToken"] = next_page_token

    resp = requests.post(endpoint, auth=auth,
                         headers=headers, json=payload, timeout=60)
    resp.raise_for_status()
    data = resp.json()

    for issue in data.get("issues", []):
        row = {"key": issue["key"]}
        for jira_fid, db_col in _JIRA_FIELD_MAP.items():
            if db_col in row:
                continue
            row[db_col] = _jira_extract(jira_fid,
                                        issue["fields"].get(jira_fid))
        all_rows.append(row)

    if progress_cb:
        progress_cb(len(all_rows), total_hint)

    if data.get("isLast", True):
        break
    next_page_token = data.get("nextPageToken")
    if not next_page_token:
        break
\end{lstlisting}

\subsection{Step 4 — Field Extraction}

The \texttt{\_jira\_extract} helper normalises the diverse value types returned by the API:

\begin{lstlisting}[style=python, caption={Field extraction helper (excerpt)}]
def _jira_extract(field_id: str, val):
    if val is None:
        return None

    # Plain objects: grab name / displayName / value
    if isinstance(val, dict):
        return (val.get("name") or val.get("displayName")
                or val.get("value") or val.get("key"))

    # Lists: join names
    if isinstance(val, list):
        parts = [str(item.get("name") or item.get("value") or "")
                 for item in val if isinstance(item, dict)]
        return ", ".join(p for p in parts if p) or None

    # ADF rich-text (Root Cause, Resolution Description, etc.)
    if isinstance(val, dict) and val.get("type") == "doc":
        return _adf_to_text(val).strip() or None

    return val  # scalar (string, number, bool)
\end{lstlisting}

\subsection{Step 5 — Transform Pipeline}

After all pages are fetched, the raw rows are assembled into a DataFrame and cleaned:

\begin{lstlisting}[style=python, caption={Transform pipeline}]
df = pd.DataFrame(all_rows)

# 1. Normalise strings
obj_cols = df.select_dtypes(include="object").columns
for col in obj_cols:
    df[col] = (df[col].astype(str).str.strip()
               .replace({"nan": None, "None": None, "": None}))

# 2. Parse all date columns (ISO strings with timezone offsets)
def _to_dt(series):
    return (pd.to_datetime(series, errors="coerce", utc=True)
            .dt.tz_convert(None))   # strip TZ -> tz-naive datetime64

for col in DATE_COLS:
    if col in df.columns:
        df[col] = _to_dt(df[col])

# 3. Merge dual story-points fields
_sp_opt = pd.to_numeric(df["story_points_option"], errors="coerce")
df["story_points"] = df["story_points"].fillna(_sp_opt)
df.drop(columns=["story_points_option"], inplace=True)

# 4. Derived columns (same as Excel ETL path)
mask = df["resolved"].notna() & df["created"].notna()
df.loc[mask, "resolution_days"] = (
    (df.loc[mask, "resolved"] - df.loc[mask, "created"])
    .dt.total_seconds() / 86400
)
df["created_yearmonth"] = df["created"].dt.to_period("M").astype(str)
df["created_year"]      = df["created"].dt.year
df["created_month"]     = df["created"].dt.month
df["created_week"]      = df["created"].dt.isocalendar().week.astype(float)
\end{lstlisting}

\subsection{Step 6 — Load into Database}

The cleaned DataFrame is written to MySQL using the existing \texttt{upsert} or
\texttt{initial\_load} functions, maintaining full schema compatibility with the Excel path:

\begin{lstlisting}[style=python, caption={Load to DB}]
from sqlalchemy import inspect as _inspect

df = df.drop_duplicates(subset=["key"], keep="last")

if mode == "replace" or not _inspect(engine).has_table("issues"):
    n = initial_load(df, engine)   # DROP + CREATE + INSERT
else:
    n = upsert(df, engine)         # DELETE existing keys + INSERT
\end{lstlisting}

% ────────────────────────────────────────────────────────────────────────────
\section{UI Integration}

A \textbf{Sync from Jira} section was added to the Upload tab of the Streamlit dashboard,
above the legacy Excel uploader. It exposes:

\begin{itemize}
    \item \textbf{Sync mode}: Incremental (issues updated since a chosen date) or
          Replace all (full reload of all $\sim$55,000 issues).
    \item \textbf{Project key}: Configurable text input, defaulting to \texttt{IRPAUTO}.
    \item \textbf{Updated since}: Date picker used in incremental mode.
    \item \textbf{Progress bar}: Live display of \texttt{fetched / total} issues as pages
          are retrieved.
\end{itemize}

\begin{lstlisting}[style=python, caption={Progress callback wired to Streamlit}]
_progress_bar = st.progress(0, text="Connecting to Jira...")

def _update_progress(done, total):
    pct = min(done / max(total, 1), 1.0)
    _progress_bar.progress(
        pct, text=f"Fetched {done:,} / {total:,} issues..."
    )

result = fetch_jira_issues(
    engine=_engine(),
    jira_url=_jira_url,
    email=_jira_email,
    api_token=_jira_token,
    project=_jira_project,
    mode=_mode_str,
    updated_since=_since_str,
    progress_cb=_update_progress,
)
\end{lstlisting}

% ────────────────────────────────────────────────────────────────────────────
\section{Technical Findings}

\subsection{API Version}

The classic \texttt{GET /rest/api/3/search?jql=...} endpoint has been deprecated and removed
by Atlassian. The replacement is \texttt{POST /rest/api/3/search/jql} which accepts a JSON
body and uses cursor-based pagination (\texttt{nextPageToken}) instead of offset-based.

\subsection{Story Points — Dual Field Issue}

IRPAUTO uses a non-standard story points configuration. Two fields exist:

\begin{itemize}
    \item \texttt{customfield\_10049} — the standard Jira story points field (numeric).
          Populated only on non-IRPAUTO issues in this instance.
    \item \texttt{customfield\_10520} — a custom dropdown field labelled ``Story Points''
          with option values \texttt{'1'}, \texttt{'2'}, \texttt{'3'}, \texttt{'4'},
          \texttt{'5'}, \texttt{'8'}, \texttt{'12'}, \texttt{'20'}, etc. This is the
          actual source of story points for IRPAUTO.
\end{itemize}

The ETL reads \texttt{customfield\_10049} first, then falls back to
\texttt{customfield\_10520}. Without this fallback, approximately 82\% of IRPAUTO issues
would have \texttt{NULL} story points, making the sprint burndown chart completely incorrect.

\subsection{Date Timezone Handling}

Jira returns all timestamps in ISO 8601 format with timezone offsets
(e.g.\ \texttt{2026-03-27T16:37:02.386+0100}). Parsing these with
\texttt{pd.to\_datetime(..., utc=True)} produces timezone-aware \texttt{datetime64[ns, UTC]}
series. The application's MySQL schema expects timezone-naive datetimes, so
\texttt{.dt.tz\_convert(None)} is applied to strip the timezone after normalisation to UTC.

\subsection{ADF Rich-Text Fields}

Fields such as Root Cause Description, Resolution Description, and Impact Description are
returned in Atlassian Document Format (ADF) — a nested JSON structure. A recursive
\texttt{\_adf\_to\_text} function traverses the document tree and extracts plain text,
preserving paragraph breaks.

\subsection{KPI SLA Fields}

The five KPI duration fields (\textit{Time to Solve}, \textit{Time Initial Response}, etc.)
are stored by Jira as complex SLA cycle objects containing start time, stop time,
breach status, and remaining time. The Excel export pre-processes these into human-readable
strings like \texttt{998h 31m}. Replicating this calculation from the raw API requires
parsing the SLA cycles, which is deferred to a future iteration. These fields are currently
stored as \texttt{NULL}.

% ────────────────────────────────────────────────────────────────────────────
\section{Recommended Workflow}

\begin{enumerate}
    \item \textbf{Initial bootstrap} (once): Run a \textit{Replace all} sync to load the
          full history of $\sim$55,000 issues. Expected duration: approximately 10 minutes.
    \item \textbf{Daily incremental sync}: Each morning, run an \textit{Incremental} sync
          with \textit{Updated since} set to yesterday's date. Only issues modified in the
          last 24 hours are fetched and upserted. Expected duration: under 1 minute for
          typical daily volumes.
    \item \textbf{After sync}: Click \textit{Refresh R25 Sprint Burndown} and
          \textit{Refresh Backlog Burndown} to recompute the pre-aggregated Power BI tables.
    \item \textbf{Release Lifecycle file}: The scope, team availability, and sprint metadata
          still require a manual upload of \texttt{Release\_lifecycle\_R25.xlsx}, as this
          data does not exist in Jira.
\end{enumerate}

% ────────────────────────────────────────────────────────────────────────────
\section{Summary}

The Jira API integration achieves full data parity with the legacy Excel extract across all
fields except the KPI SLA durations. The incremental sync mode makes day-to-day operation
fast and lightweight, while the full reload option provides a clean baseline whenever needed.
The dual story-points field handling is a critical IRPAUTO-specific detail that ensures
burndown charts and sprint metrics remain accurate.

\vspace{1cm}
\noindent\rule{\linewidth}{0.4pt}
{\small\textit{Document generated on April 2, 2026 — DXC Technology, IRP AUTO Analytics Platform.}}

\end{document}
